% !PRINTED 451
\weberpage{451}
%<<CONT
lichen (mit Einschluss der Constanten) dienen, wobei man die constanten Co\"efficienten auf einen beliebigen Zahlk\"orper, etwa auf den der rationalen Zahlen beschr\"anken oder auch ganz unbeschr\"ankt annehmen kann.

Wir wollen hier nicht weiter die Beispiele h\"aufen, da die genaue Untersuchung von K\"orpern besonderer Art eine der Hauptaufgaben unserer sp\"ateren Ausf\"uhrungen sein wird.

\begin{center}
\S.\,140.

\gesperrt{Adjunction.}
\end{center}

Wenn einem K\"orper irgend welcher Gr\"ossen, den wir mit $\Omega$ bezeichnen wollen, irgend eine Gr\"osse $\alpha$ hinzugef\"ugt wird, die nicht in ihm enthalten ist, so entsteht ein neues Gr\"ossensystem
\[
\Omega,\ \alpha,
\]
das aber kein K\"orper sein wird; um daraus einen K\"orper abzuleiten, muss man alle Gr\"ossen hinzuf\"ugen, die sich durch die Verbindung von $\alpha$ mit den Zahlen von $\Omega$ mittelst der Grundrechnungsarten ableiten lassen. So erh\"alt man einen erweiterten K\"orper $\Omega'$, der zugleich $\alpha$ und $\Omega$ enth\"alt, und der durch $\alpha$ und $\Omega$ v\"ollig bestimmt ist. Wir wollen ihn den K\"orper $\Omega, \alpha$ nennen.

Dies Hinzuf\"ugen einer neuen Gr\"osse $\alpha$ zu einem K\"orper $\Omega$ heisst \gesperrt{Adjungiren}, und man sagt, der K\"orper $\Omega'$ entsteht aus $\Omega$ durch \gesperrt{Adjunction von} $\alpha$.

Zu $\Omega'$ kann man wieder eine Gr\"osse $\alpha'$ adjungiren, und erh\"alt einen dritten K\"orper $\Omega''$, der dann $\Omega$, $\alpha$ und $\alpha'$ enth\"alt u.\,s.\,f. Man kann aber gleichzeitig zwei oder mehrere Zahlen zu $\Omega$ adjungiren, und erh\"alt so denselben K\"orper $\Omega''$, wenn man gleichzeitig $\alpha$ und $\alpha'$ dem $\Omega$ adjungirt. Man kann auch einem K\"orper einen zweiten K\"orper adjungiren und erh\"alt so einen K\"orper, der zugleich die Zahlen der beiden K\"orper enth\"alt. In dem obigen Beispiel w\"urde man durch Adjunction des K\"orpers $\Omega, \alpha'$ zu dem K\"orper $\Omega' = \Omega, \alpha$ wieder denselben K\"orper $\Omega''$ erhalten.

Adjungirt man einem K\"orper einen seiner Theile, so entsteht kein neuer K\"orper.

So erh\"alt man z.\,B., wenn man zum K\"orper der rationalen Zahlen, der mit $\mathfrak{R}$ bezeichnet sein mag, die Zahl $i=\sqrt{-1}$
%CONT>>
